\documentclass[11pt,a4paper]{article}
% Shared preamble for RuPhO English problem statements (match T1 2026 style).
% Use: \input{latex/rupho_en_preamble.tex} after \documentclass.
\usepackage[margin=1in]{geometry}
\usepackage{amsmath,amssymb}
\usepackage{graphicx}
\usepackage{float}
\usepackage{enumitem}
\usepackage{microtype}
\usepackage{caption}
\usepackage{tabularx}
\usepackage{hyperref}

\captionsetup{font=small,labelfont=bf,skip=6pt}

\setlist[itemize]{leftmargin=1.2cm}
\setlist[enumerate]{leftmargin=1.2cm}

% One outer box per subproblem: [id | statement | points] (no inner rules)
\newcommand{\problembox}[3]{%
  \par\addvspace{0.42em}%
  \noindent
  \setlength{\fboxsep}{7.5pt}%
  \setlength{\fboxrule}{0.95pt}%
  \fbox{%
    \begin{minipage}{\dimexpr\linewidth-2\fboxsep-2\fboxrule\relax}
    \begin{tabularx}{\linewidth}{@{}>{\raggedright\arraybackslash\bfseries}p{2.1em} @{\hspace{0.9em}} >{\raggedright\arraybackslash}X @{\hspace{0.9em}} >{\raggedleft\arraybackslash\bfseries}p{2.4em}@{}}
    #1 & #2 & #3
    \end{tabularx}
    \end{minipage}%
  }%
  \par\addvspace{0.32em}%
}

\title{\textbf{Steam in a test tube}}
\author{\normalsize Y17 (2017) --- English translation}
\date{}
\begin{document}
\maketitle

A thin test tube is partially filled with water and placed vertically in a large volume chamber in which the air is maintained at zero relative humidity and a constant pressure equal to external atmospheric pressure. Due to diffusion in the test tube, a linear change in the vapor concentration with height is established: near the surface of the water, the vapor turns out to be saturated, and at the upper open end of the test tube its concentration is $2$ times less. The top of the test tube is closed with a lid. After equilibrium was established, it turned out that the density of moist air inside the test tube differs from the density of dry air away from the test tube by $\Delta{\rho}=5.4~\text{g}/\text{m}^3$. Molar mass of dry air $\mu_\text{in}=29~\text{g}/\text{mol}$, vapor $\mu_\text{p}=18~\text{g}/\text{mol}$. Neglect the change in the liquid level in the test tube during the experiment.

\problembox{A1}{Find the saturated vapor pressure at the experimental temperature equal to $T=27^{\circ}\text{C}$.}{3.00}

\vspace{1em}
\noindent\textit{Source:} \url{https://pho.rs/p/3036}

\end{document}